\chapter{Nombres rationnels et opérations}

\noindent\textit{Tous les nombres que tu manipules depuis le primaire — entiers, décimaux,
fractions — appartiennent à un même grand ensemble : les \textbf{nombres rationnels}. Dans ce
chapitre, tu vas apprendre à écrire une fraction sous sa forme la plus simple, à comparer deux
rationnels, et à mener sans erreur une suite d'opérations en respectant les \textbf{priorités
opératoires}. Ces techniques sont la base de tout le calcul en classe de 3\ieme{}.}
\vspace{8pt}

\connecteBanner

\section{Entiers, décimaux, rationnels}

\begin{definition}[Nombre rationnel]
Un \textbf{nombre rationnel} est un nombre qui peut s'écrire sous la forme d'une fraction
$\dfrac{a}{b}$, où $a$ et $b$ sont des \textbf{entiers relatifs} et $b\ne0$. L'ensemble des
nombres rationnels se note $\Q$.
\end{definition}

\begin{remarque}
Tout entier est rationnel ($5=\dfrac51$), tout nombre décimal est rationnel
($0{,}75=\dfrac{75}{100}$). On a donc l'inclusion $\N\subset\Z\subset\Q$. Certains nombres,
comme $\sqrt2$ ou $\pi$, ne sont \textbf{pas} rationnels : on les étudiera plus loin.
\end{remarque}

\begin{exemple}
$-\dfrac37$, $4$, $-1{,}25=-\dfrac54$ et $\dfrac{22}{7}$ sont des rationnels. En revanche
$\sqrt2\approx1{,}414\ldots$ n'est \textbf{pas} un rationnel (ses décimales ne se répètent
jamais).
\end{exemple}

\section{Écriture fractionnaire et simplification}

\begin{definition}[Fractions égales]
Deux fractions $\dfrac ab$ et $\dfrac cd$ (avec $b,d\ne0$) sont \textbf{égales} lorsque
$a\times d=b\times c$ (produit en croix). On dit aussi qu'elles représentent le même
rationnel.
\end{definition}

\begin{propriete}[Simplification d'une fraction]
On ne change pas la valeur d'une fraction en multipliant \textbf{ou} en divisant son
numérateur et son dénominateur par un même nombre non nul. Une fraction est
\textbf{irréductible} lorsque son numérateur et son dénominateur n'ont plus de diviseur
commun (autre que $1$).
\end{propriete}

\begin{exemple}
$\dfrac{45}{75}$ : $45$ et $75$ sont tous deux divisibles par $15$, donc
$\dfrac{45}{75}=\dfrac{45\div15}{75\div15}=\dfrac35$. La fraction $\dfrac35$ est irréductible.
\end{exemple}

\begin{illus}
\begin{tikzpicture}[scale=1]
\draw[fill=onbsoft,draw=onbline,line width=0.8pt] (0,0) rectangle (6,1);
\foreach \i in {1,...,5} \draw[onbline] (\i,0) -- (\i,1);
\draw[fill=onbo] (0,0) rectangle (4,1);
\node[below] at (3,-0.15) {\small $\dfrac{4}{6}$ (6 parts, 4 coloriées)};
\draw[fill=onbsoft,draw=onbline,line width=0.8pt] (8,0) rectangle (11,1);
\draw[onbline] (9,0) -- (9,1);
\draw[fill=onbo] (8,0) rectangle (10,1);
\node[below] at (9.5,-0.15) {\small $\dfrac{2}{3}$ (3 parts, 2 coloriées)};
\draw[onbmuted,<->,>=Stealth] (6.3,0.5) -- (7.7,0.5);
\node[above,onbmuted] at (7,0.7) {\small même aire};
\end{tikzpicture}
\fig{Figure — Deux fractions égales : $\dfrac46=\dfrac23$, la même surface coloriée
partagée différemment.}
\end{illus}

\begin{aretenir}
Pour simplifier une fraction, on peut décomposer numérateur et dénominateur en produits de
petits facteurs communs, ou diviser plusieurs fois de suite par un diviseur commun évident
(2, 3, 5, ...), jusqu'à ce qu'il n'y en ait plus.
\end{aretenir}

\section{Comparaison et ordre}

\begin{propriete}[Comparer deux rationnels]
Pour comparer deux fractions, on les réduit \textbf{au même dénominateur} (positif), puis on
compare leurs numérateurs. On peut aussi comparer $\dfrac ab$ et $\dfrac cd$ ($b,d>0$) en
comparant les produits en croix $a\times d$ et $b\times c$.
\end{propriete}

\begin{exemple}
Comparer $\dfrac56$ et $\dfrac79$ : dénominateur commun $18$. $\dfrac56=\dfrac{15}{18}$ et
$\dfrac79=\dfrac{14}{18}$. Comme $15>14$, on a $\dfrac56>\dfrac79$.
\end{exemple}

\begin{illus}
\begin{tikzpicture}[>=Stealth]
\draw[onbline,line width=1pt] (-0.3,0) -- (6.3,0);
\foreach \x/\lbl in {0/-1,1.5/-\tfrac12,3/0,4.5/\tfrac12,6/1}
  \draw[onbmuted] (\x,0.08) -- (\x,-0.08) node[below] {\small $\lbl$};
\filldraw[onbo] (3.75,0) circle (2.2pt) node[above] {\small $\tfrac14$};
\filldraw[onbgreen] (2.4,0) circle (2.2pt) node[above] {\small $-\tfrac15$};
\filldraw[onbblue] (5.4,0) circle (2.2pt) node[above] {\small $\tfrac45$};
\end{tikzpicture}
\fig{Figure — Droite graduée : $-\tfrac15<\tfrac14<\tfrac45$, l'ordre se lit de gauche à
droite.}
\end{illus}

\begin{definition}[Valeur absolue]
La \textbf{valeur absolue} d'un rationnel $x$, notée $|x|$, est sa distance à $0$ sur la
droite graduée. On a $|x|=x$ si $x\ge0$ et $|x|=-x$ si $x<0$ ; $|x|$ est toujours positif ou
nul.
\end{definition}

\begin{exemple}
$\left|-\dfrac74\right|=\dfrac74$ et $\left|\dfrac53\right|=\dfrac53$. Comme
$\dfrac74=1{,}75>\dfrac53\approx1{,}67$, le nombre $-\dfrac74$ est plus \emph{éloigné} de $0$
que $\dfrac53$.
\end{exemple}

\begin{remarque}
Entre deux nombres négatifs, celui qui a la plus \textbf{grande} valeur absolue est le plus
\textbf{petit} : $-\dfrac74<-\dfrac53$ bien que $\dfrac74>\dfrac53$.
\end{remarque}

\section{Addition et soustraction de rationnels}

\begin{propriete}[Somme et différence]
Pour additionner (ou soustraire) deux fractions, on les réduit d'abord au \textbf{même
dénominateur}, puis on additionne (ou soustrait) les numérateurs en gardant le dénominateur
commun : $\dfrac ab+\dfrac cb=\dfrac{a+c}b$.
\end{propriete}

\begin{illus}
\begin{tikzpicture}
\draw[fill=onbsoft,draw=onbline,line width=0.8pt] (0,0) rectangle (4,1);
\foreach \i in {1,2,3} \draw[onbline] (\i,0) -- (\i,1);
\draw[fill=onbo] (0,0) rectangle (1,1);
\draw[fill=onbo2] (1,0) rectangle (3,1);
\node[below] at (2,-0.2) {\small $\dfrac14+\dfrac24=\dfrac34$};
\end{tikzpicture}
\fig{Figure — Addition de fractions de même dénominateur : les parts coloriées
s'ajoutent bout à bout.}
\end{illus}

\begin{exemple}
$\dfrac23+\dfrac54-\dfrac16$ : dénominateur commun $12$. On obtient
$\dfrac8{12}+\dfrac{15}{12}-\dfrac2{12}=\dfrac{8+15-2}{12}=\dfrac{21}{12}=\dfrac74$.
\end{exemple}

\begin{aretenir}
Ne jamais additionner (ni soustraire) des numérateurs de fractions qui n'ont pas le
\textbf{même} dénominateur : $\dfrac12+\dfrac13\ne\dfrac25$ ! Il faut toujours passer par un
dénominateur commun (le produit des dénominateurs convient toujours, le PPCM donne des
nombres plus petits).
\end{aretenir}

\section{Multiplication et division de rationnels}

\begin{propriete}[Produit et quotient]
$\dfrac ab\times\dfrac cd=\dfrac{a\times c}{b\times d}$. Diviser par une fraction non nulle
revient à multiplier par son \textbf{inverse} :
$\dfrac ab\div\dfrac cd=\dfrac ab\times\dfrac dc$ (pour $c\ne0$). On simplifie si possible
\textbf{avant} de multiplier, cela évite de gros calculs.
\end{propriete}

\begin{exemple}
$\dfrac34\times\left(-\dfrac89\right)=-\dfrac{3\times8}{4\times9}=-\dfrac{24}{36}=-\dfrac23$
(en simplifiant par $12$). \quad
$\left(-\dfrac56\right)\div\dfrac{10}3=-\dfrac56\times\dfrac3{10}=-\dfrac{15}{60}=-\dfrac14$.
\end{exemple}

\begin{remarque}
La division par $0$ est \textbf{impossible} : $\dfrac c0$ n'a aucun sens, quelle que soit la
valeur de $c$. Le signe d'un produit ou d'un quotient suit la règle des signes habituelle :
$(-)\times(-)=(+)$, $(+)\times(-)=(-)$, etc.
\end{remarque}

\section{Priorités opératoires}

\begin{aretenir}[Ordre des priorités]
Dans un calcul sans parenthèses, on effectue \textbf{dans l'ordre} : \textbf{1)} les
puissances (vues au chapitre suivant), \textbf{2)} les multiplications et divisions, de
gauche à droite, \textbf{3)} les additions et soustractions, de gauche à droite. S'il y a des
\textbf{parenthèses}, on calcule d'abord ce qu'elles contiennent, en commençant par les plus
intérieures.
\end{aretenir}

\begin{illus}
\begin{tikzpicture}[node distance=10mm,>=Stealth,every node/.style={font=\small}]
\node[draw=onbo,fill=onbsoft,rounded corners=3pt,minimum width=2.6cm,minimum height=0.8cm] (p1) {1. Parenthèses};
\node[draw=onbblue,fill=white,rounded corners=3pt,minimum width=2.6cm,minimum height=0.8cm,right=of p1] (p2) {2. $\times\;\div$};
\node[draw=onbgreen,fill=white,rounded corners=3pt,minimum width=2.6cm,minimum height=0.8cm,right=of p2] (p3) {3. $+\;-$};
\draw[->,onbmuted,line width=1pt] (p1) -- (p2);
\draw[->,onbmuted,line width=1pt] (p2) -- (p3);
\end{tikzpicture}
\fig{Figure — L'ordre des priorités opératoires : parenthèses, puis multiplication/division,
puis addition/soustraction.}
\end{illus}

\begin{exemple}
$5-2\times\left(\dfrac34-\dfrac12\right)+\dfrac18$ : on calcule d'abord la parenthèse
$\dfrac34-\dfrac12=\dfrac14$, puis la multiplication $2\times\dfrac14=\dfrac12$, et enfin
$5-\dfrac12+\dfrac18=\dfrac{40}8-\dfrac48+\dfrac18=\dfrac{37}8$.
\end{exemple}

% ═══════════════════════════════════════════════════════════════
\section{L'essentiel à retenir}

\begin{synthese}
\textbf{Rationnels.} $\Q$ = nombres $\dfrac ab$ avec $a,b\in\Z$, $b\ne0$. $\N\subset\Z\subset\Q$.\\[4pt]
\textbf{Fractions égales.} $\dfrac ab=\dfrac cd\iff a\times d=b\times c$ (produit en croix).\\[4pt]
\textbf{Simplifier.} Diviser numérateur et dénominateur par un même diviseur commun, jusqu'à
la forme irréductible.\\[4pt]
\textbf{Comparer.} Même dénominateur (positif) $\Rightarrow$ comparer les numérateurs.
Valeur absolue $|x|$ = distance de $x$ à $0$.\\[4pt]
\textbf{Addition/soustraction.} $\dfrac ab\pm\dfrac cb=\dfrac{a\pm c}b$ (même dénominateur
obligatoire).\\[4pt]
\textbf{Multiplication/division.} $\dfrac ab\times\dfrac cd=\dfrac{ac}{bd}$ ;\ diviser =
multiplier par l'inverse. Division par $0$ interdite.\\[4pt]
\textbf{Priorités.} Parenthèses $\to$ $\times,\div$ (gauche à droite) $\to$ $+,-$ (gauche à
droite).\\[4pt]
\textbf{Piège fréquent.} $\dfrac12+\dfrac13\ne\dfrac25$ : jamais additionner des fractions de
dénominateurs différents sans les réduire au même dénominateur d'abord.
\end{synthese}

% ═══════════════════════════════════════════════════════════════
\section{Méthodes \& savoir-faire}

\methode{Simplifier une fraction jusqu'à sa forme irréductible}
Chercher un diviseur commun évident au numérateur et au dénominateur (2, 3, 5, ...), diviser,
et recommencer jusqu'à ce qu'il n'y en ait plus.
\begin{exemple}
$\dfrac{45}{75}$ : diviser par $5$ donne $\dfrac9{15}$, puis diviser par $3$ donne
$\dfrac35$, qui est irréductible.
\end{exemple}

\methode{Comparer deux fractions}
Réduire les deux fractions au même dénominateur (positif) puis comparer les numérateurs.
\begin{exemple}
$-\dfrac56$ et $-\dfrac79$ : dénominateur commun $18$ donne $-\dfrac{15}{18}$ et
$-\dfrac{14}{18}$. Comme $-15<-14$, $-\dfrac56<-\dfrac79$.
\end{exemple}

\methode{Additionner ou soustraire des fractions de dénominateurs différents}
Trouver un dénominateur commun (le produit des dénominateurs, ou leur PPCM), réécrire chaque
fraction, puis additionner/soustraire les numérateurs.
\begin{exemple}
$\dfrac34-\dfrac56$ : dénominateur commun $12$, donne $\dfrac9{12}-\dfrac{10}{12}=-\dfrac1{12}$.
\end{exemple}

\methode{Multiplier et diviser des fractions}
Multiplier numérateurs entre eux et dénominateurs entre eux (en simplifiant si possible
avant) ; pour diviser, multiplier par l'inverse du diviseur.
\begin{exemple}
$\dfrac23\times\left(-\dfrac9{10}\right)=-\dfrac{2\times9}{3\times10}=-\dfrac{18}{30}=-\dfrac35$
(simplifié par $6$). \quad $\dfrac78\div\left(-\dfrac{14}3\right)=\dfrac78\times\left(-\dfrac3{14}\right)=-\dfrac{21}{112}=-\dfrac3{16}$.
\end{exemple}

\methode{Appliquer correctement les priorités opératoires}
Repérer d'abord les parenthèses, les calculer entièrement, puis effectuer les multiplications
et divisions avant les additions et soustractions.
\begin{exemple}
$4-3\times\left(\dfrac23-1\right)$ : parenthèse $\dfrac23-1=-\dfrac13$, puis
$3\times\left(-\dfrac13\right)=-1$, puis $4-(-1)=5$.
\end{exemple}

\methode{Utiliser la valeur absolue pour comparer des rationnels}
Calculer les valeurs absolues des deux nombres : si les deux nombres sont négatifs, celui qui
a la plus grande valeur absolue est le plus petit.
\begin{exemple}
$\left|-\dfrac74\right|=\dfrac74=1{,}75$ et $\left|\dfrac53\right|=\dfrac53\approx1{,}67$.
Comme $1{,}75>1{,}67$, on a $-\dfrac74<\dfrac53$ (évident aussi car le premier est négatif).
\end{exemple}

\methode{Résoudre un problème concret avec des rationnels}
Traduire chaque étape de l'énoncé par une opération sur des fractions, en calculant chaque
quantité intermédiaire avant de passer à la suivante.
\begin{exemple}
Kouam a $54\,000$ FCFA. Il donne $\dfrac16$ de cette somme à son frère, puis dépense
$\dfrac29$ du reste pour des habits. Il donne : $\dfrac16\times54\,000=9\,000$ FCFA (reste
$45\,000$ FCFA) ; habits : $\dfrac29\times45\,000=10\,000$ FCFA. Il lui reste
$45\,000-10\,000=35\,000$ FCFA.
\end{exemple}

% ═══════════════════════════════════════════════════════════════
\section{Exercices d'application}
\begin{multicols}{2}

\rubrique{Simplifier et comparer des fractions}
\exo{}\diff{1} Simplifier jusqu'à la forme irréductible : $\dfrac{36}{48}$,\ $\dfrac{56}{84}$,\ $\dfrac{40}{64}$.

\exo{}\diff{1} Simplifier jusqu'à la forme irréductible : $-\dfrac{84}{126}$,\ $-\dfrac{56}{98}$.

\exo{}\diff{2} Comparer $\dfrac56$ et $\dfrac79$ (méthode du dénominateur commun).

\exo{}\diff{2} Ranger dans l'ordre croissant : $-\dfrac34$,\ $-\dfrac23$,\ $-\dfrac58$.

\rubrique{Addition, soustraction, multiplication, division}
\exo{}\diff{1} Calculer $A=\dfrac23+\dfrac54-\dfrac16$ (résultat sous forme de fraction
irréductible).

\exo{}\diff{1} Calculer $B=-\dfrac7{10}-\dfrac35+\dfrac{11}{20}$.

\exo{}\diff{2} Calculer $C=\dfrac34\times\left(-\dfrac89\right)$.

\exo{}\diff{2} Calculer $D=\left(-\dfrac56\right)\div\dfrac{10}3$.

\rubrique{Priorités opératoires}
\exo{}\diff{2} Calculer $E=\dfrac25+\dfrac34\times\left(-\dfrac23\right)$.

\exo{}\diff{2} Calculer $F=\left(\dfrac12-\dfrac23\right)\times\left(\dfrac34+\dfrac16\right)$.

\exo{}\diff{2} Calculer $G=5-2\times\left(\dfrac34-\dfrac12\right)+\dfrac18$.

\exo{}\diff{3} Calculer $H=\left[2-3\times\left(\dfrac12-\dfrac56\right)\right]\div\left(-\dfrac23\right)$.

\rubrique{Problèmes concrets}
\exo{}\diff{2} Un champ de $3\,600\text{ m}^2$, situé près de Bafoussam, est partagé entre
trois héritiers. Le premier reçoit $\dfrac13$ de la superficie totale, le deuxième reçoit
$\dfrac25$ de ce qui reste, et le troisième reçoit le reste. Calculer la superficie reçue par
chaque héritier.

\exo{}\diff{2} Mireille dispose de $45\,000$ FCFA. Elle dépense $\dfrac25$ de cette somme en
nourriture, puis $\dfrac13$ de ce qui reste en transport. Combien d'argent lui reste-t-il ?

\exo{}\diff{2} Un taxi doit parcourir un trajet de $96$ km entre Douala et une ville voisine.
Le matin, il parcourt $\dfrac38$ du trajet ; l'après-midi, $\dfrac5{12}$ du trajet. Quelle
distance (en km) lui reste-t-il à parcourir ?

\exo{}\diff{1} Une recette pour $8$ personnes nécessite $\dfrac34$ kg de farine. En gardant
les proportions, quelle quantité de farine faut-il pour $6$ personnes ?

% ═══════════════════════════════════════════════════════════════
\end{multicols}

\section{Exercices d'entraînement}
\begin{multicols}{2}

\rubrique{Simplifier et comparer}
\exo{}\diff{1} Simplifier : $\dfrac{28}{42}$,\ $\dfrac{63}{99}$,\ $\dfrac{48}{60}$.

\exo{}\diff{1} Simplifier : $-\dfrac{72}{108}$,\ $-\dfrac{90}{135}$.

\exo{}\diff{2} Comparer $\dfrac{11}{15}$ et $\dfrac{13}{18}$.

\exo{}\diff{2} Comparer $-\dfrac45$ et $-\dfrac79$.

\exo{}\diff{2} Ranger dans l'ordre décroissant : $\dfrac53$,\ $\dfrac74$,\ $\dfrac32$.

\rubrique{Opérations sur les rationnels}
\exo{}\diff{1} Calculer $\dfrac37+\dfrac27-\dfrac57$.

\exo{}\diff{1} Calculer $-\dfrac54+\dfrac32$.

\exo{}\diff{2} Calculer $\dfrac59-\dfrac13+\dfrac16$.

\exo{}\diff{2} Calculer $\left(-\dfrac47\right)\times\dfrac{14}5$.

\exo{}\diff{2} Calculer $\dfrac{9}{10}\div\left(-\dfrac35\right)$.

\rubrique{Priorités opératoires}
\exo{}\diff{2} Calculer $\dfrac13+\dfrac25\times\dfrac52$.

\exo{}\diff{2} Calculer $\left(3-\dfrac54\right)\times\left(2+\dfrac13\right)$.

\exo{}\diff{2} Calculer $7-3\times\left(\dfrac23-\dfrac14\right)$.

\exo{}\diff{3} Calculer $\left[\dfrac54-\left(2-\dfrac73\right)\right]\times\dfrac{-6}5$.

\exo{}\diff{3} Calculer $\dfrac{4-\dfrac23\times3}{1-\dfrac12}$.

\rubrique{Valeur absolue et ordre}
\exo{}\diff{1} Calculer $\left|-\dfrac97\right|$ et $\left|\dfrac{-11}4\right|$.

\exo{}\diff{2} Sans calculer leur valeur exacte, comparer $-\dfrac{17}5$ et $-\dfrac{16}5$ à
l'aide de leur valeur absolue.

\exo{}\diff{2} Placer approximativement sur une droite graduée les nombres $-\dfrac32$,
$\dfrac14$, $-\dfrac25$, $\dfrac95$.

\exo{}\diff{2} Trouver un rationnel compris strictement entre $\dfrac23$ et $\dfrac34$.

\rubrique{Problèmes concrets}
\exo{}\diff{2} Un réservoir contient $2\,400$ litres d'eau. On en utilise $\dfrac38$ pour
arroser un champ à Maroua, puis $\dfrac14$ du reste pour le bétail. Quelle quantité d'eau
reste-t-il dans le réservoir ?

\exo{}\diff{2} Une classe de 3\ieme{} compte $60$ élèves. Les $\dfrac25$ des élèves pratiquent
le football, et $\dfrac13$ des élèves restants pratiquent le basket. Combien d'élèves ne
pratiquent ni l'un ni l'autre de ces deux sports (parmi les élèves considérés) ?

\exo{}\diff{2} Un commerçant de Kribi achète des ananas pour $84\,000$ FCFA. Il revend les
$\dfrac57$ de sa marchandise avec un bénéfice, puis constate qu'il lui reste encore
$\dfrac13$ du reste à vendre le lendemain. Quelle fraction de la marchandise totale a-t-il
vendue le premier jour, et quelle fraction lui reste-t-il après le début du deuxième jour ?

\exo{}\diff{3} Un terrain rectangulaire de $5\,000\text{ m}^2$ à Ebolowa est partagé en trois
lots : le premier occupe $\dfrac25$ du terrain, le deuxième occupe $\dfrac13$ du reste, et le
troisième lot correspond à ce qui reste. Calculer la superficie de chaque lot.

\exo{}\diff{2} Un litre de jus de fruit coûte $650$ FCFA. Adèle achète $\dfrac34$ de litre.
Quel est le prix payé ?

% ═══════════════════════════════════════════════════════════════
\end{multicols}

\section{Sujets type BEPC}

\sujetbac[fractions, comparaison et partage — marché du Mfoundi, Yaoundé]
\begin{enumerate}[label=\arabic*)]
\item Simplifier la fraction $\dfrac{72}{120}$ jusqu'à sa forme irréductible.
\item Comparer les rationnels $-\dfrac56$ et $-\dfrac79$.
\item Calculer et donner le résultat sous forme de fraction irréductible :
$\left(\dfrac23-\dfrac14\right)\times\left(-\dfrac35\right)+\dfrac12$.
\item Un commerçant du marché du Mfoundi, à Yaoundé, achète un sac d'arachides de $80$ kg
pour $32\,000$ FCFA. Il en revend les $\dfrac35$ au détail à $500$ FCFA le kilogramme, et le
reste à un grossiste à $350$ FCFA le kilogramme.
\begin{enumerate}[label=\alph*)]
\item Calculer la masse (en kg) vendue au détail, puis la masse vendue au grossiste.
\item Calculer la recette totale du commerçant.
\item Calculer son bénéfice (recette totale moins prix d'achat).
\end{enumerate}
\end{enumerate}

\sujetbac[fractions et trajet — gare routière de Douala]
\begin{enumerate}[label=\arabic*)]
\item Simplifier la fraction $\dfrac{105}{135}$ jusqu'à sa forme irréductible.
\item Comparer les rationnels $\dfrac{11}{15}$ et $\dfrac{13}{18}$.
\item Calculer $D=3-\left(\dfrac23+\dfrac16\right)\times\dfrac45$.
\item Un bus interurbain relie Douala à Bafoussam, sur une distance totale de $300$ km. Il
parcourt $\dfrac25$ du trajet avant une première pause, puis $\dfrac7{15}$ du trajet avant une
seconde pause.
\begin{enumerate}[label=\alph*)]
\item Calculer, en km, la distance parcourue avant chaque pause.
\item Quelle fraction du trajet reste-t-il à parcourir après la seconde pause ? Donner la
distance correspondante en km.
\end{enumerate}
\end{enumerate}

\sujetbac[fractions et économies — Garoua]
\begin{enumerate}[label=\arabic*)]
\item Simplifier la fraction $\dfrac{130}{195}$ jusqu'à sa forme irréductible.
\item Comparer les rationnels $-\dfrac45$ et $-\dfrac79$.
\item Calculer $E=\left(\dfrac54-\dfrac13\right)\times2-\dfrac38$.
\item Aïcha, à Garoua, dispose de $60\,000$ FCFA d'économies. Elle place $\dfrac25$ de cette
somme à la banque, dépense ensuite $\dfrac14$ de ce qui reste pour des fournitures scolaires,
puis donne le tiers de la nouvelle somme restante à sa petite sœur.
\begin{enumerate}[label=\alph*)]
\item Calculer la somme placée à la banque, puis la somme dépensée en fournitures.
\item Calculer la somme donnée à sa petite sœur.
\item Combien d'argent reste-t-il finalement à Aïcha ?
\end{enumerate}
\end{enumerate}

% ═══════════════════════════════════════════════════════════════
\section{Corrigés}
\begin{multicols}{2}

\corrige{de l'exercice 1}
$\dfrac{36}{48}=\dfrac{36\div12}{48\div12}=\dfrac34$ ;\quad
$\dfrac{56}{84}=\dfrac{56\div28}{84\div28}=\dfrac23$ ;\quad
$\dfrac{40}{64}=\dfrac{40\div8}{64\div8}=\dfrac58$.

\corrige{de l'exercice 2}
$-\dfrac{84}{126}=-\dfrac{84\div42}{126\div42}=-\dfrac23$ ;\quad
$-\dfrac{56}{98}=-\dfrac{56\div14}{98\div14}=-\dfrac47$.

\corrige{de l'exercice 3}
Dénominateur commun $18$ : $\dfrac56=\dfrac{15}{18}$ et $\dfrac79=\dfrac{14}{18}$. Comme
$15>14$, $\dfrac56>\dfrac79$.

\corrige{de l'exercice 4}
Dénominateur commun $24$ : $-\dfrac34=-\dfrac{18}{24}$, $-\dfrac23=-\dfrac{16}{24}$,
$-\dfrac58=-\dfrac{15}{24}$. Ordre croissant (numérateurs négatifs croissants) :
$-\dfrac34<-\dfrac23<-\dfrac58$.

\corrige{de l'exercice 5}
Dénominateur commun $12$ : $A=\dfrac8{12}+\dfrac{15}{12}-\dfrac2{12}=\dfrac{21}{12}=\dfrac74$.

\corrige{de l'exercice 6}
Dénominateur commun $20$ : $B=-\dfrac{14}{20}-\dfrac{12}{20}+\dfrac{11}{20}=-\dfrac{15}{20}=-\dfrac34$.

\corrige{de l'exercice 7}
$C=\dfrac{3\times(-8)}{4\times9}=-\dfrac{24}{36}=-\dfrac23$.

\corrige{de l'exercice 8}
$D=-\dfrac56\times\dfrac3{10}=-\dfrac{15}{60}=-\dfrac14$.

\corrige{de l'exercice 9}
$E=\dfrac25+\dfrac34\times\left(-\dfrac23\right)=\dfrac25-\dfrac12=\dfrac4{10}-\dfrac5{10}=-\dfrac1{10}$.

\corrige{de l'exercice 10}
$\dfrac12-\dfrac23=-\dfrac16$ ;\ $\dfrac34+\dfrac16=\dfrac{11}{12}$. Donc
$F=-\dfrac16\times\dfrac{11}{12}=-\dfrac{11}{72}$.

\corrige{de l'exercice 11}
Parenthèse : $\dfrac34-\dfrac12=\dfrac14$. Puis $2\times\dfrac14=\dfrac12$. Enfin
$G=5-\dfrac12+\dfrac18=\dfrac{40}8-\dfrac48+\dfrac18=\dfrac{37}8$.

\corrige{de l'exercice 12}
Parenthèse : $\dfrac12-\dfrac56=-\dfrac13$. Puis $3\times\left(-\dfrac13\right)=-1$, donc
$2-(-1)=3$. Enfin $H=3\div\left(-\dfrac23\right)=3\times\left(-\dfrac32\right)=-\dfrac92$.

\corrige{de l'exercice 13}
Premier héritier : $\dfrac13\times3\,600=1\,200\text{ m}^2$ ; il reste
$3\,600-1\,200=2\,400\text{ m}^2$. Deuxième héritier : $\dfrac25\times2\,400=960\text{ m}^2$.
Troisième héritier : $2\,400-960=1\,440\text{ m}^2$.

\corrige{de l'exercice 14}
Nourriture : $\dfrac25\times45\,000=18\,000$ FCFA ; il reste $45\,000-18\,000=27\,000$ FCFA.
Transport : $\dfrac13\times27\,000=9\,000$ FCFA. Il lui reste $27\,000-9\,000=18\,000$ FCFA.

\corrige{de l'exercice 15}
Fraction déjà parcourue : $\dfrac38+\dfrac5{12}=\dfrac9{24}+\dfrac{10}{24}=\dfrac{19}{24}$.
Fraction restante : $1-\dfrac{19}{24}=\dfrac5{24}$. Distance restante :
$\dfrac5{24}\times96=20$ km.

\corrige{de l'exercice 16}
Le facteur d'échelle est $\dfrac68=\dfrac34$. Farine nécessaire :
$\dfrac34\times\dfrac34=\dfrac9{16}$ kg.

\corrige{du sujet 1}
1) $\dfrac{72}{120}=\dfrac{72\div24}{120\div24}=\dfrac35$.\\
2) Dénominateur commun $18$ : $-\dfrac56=-\dfrac{15}{18}$, $-\dfrac79=-\dfrac{14}{18}$. Comme
$-15<-14$, $-\dfrac56<-\dfrac79$.\\
3) $\dfrac23-\dfrac14=\dfrac8{12}-\dfrac3{12}=\dfrac5{12}$ ; puis
$\dfrac5{12}\times\left(-\dfrac35\right)=-\dfrac{15}{60}=-\dfrac14$ ; enfin
$-\dfrac14+\dfrac12=\dfrac14$.\\
4a) Détail : $\dfrac35\times80=48$ kg ; grossiste : $80-48=32$ kg.\\
4b) Recette détail : $48\times500=24\,000$ FCFA ; recette grossiste : $32\times350=11\,200$
FCFA ; recette totale : $24\,000+11\,200=35\,200$ FCFA.\\
4c) Bénéfice : $35\,200-32\,000=3\,200$ FCFA.

\corrige{du sujet 2}
1) $\dfrac{105}{135}=\dfrac{105\div15}{135\div15}=\dfrac79$.\\
2) Dénominateur commun $90$ : $\dfrac{11}{15}=\dfrac{66}{90}$, $\dfrac{13}{18}=\dfrac{65}{90}$.
Comme $66>65$, $\dfrac{11}{15}>\dfrac{13}{18}$.\\
3) $\dfrac23+\dfrac16=\dfrac56$ ; puis $\dfrac56\times\dfrac45=\dfrac{20}{30}=\dfrac23$ ; enfin
$D=3-\dfrac23=\dfrac73$.\\
4a) Avant la première pause : $\dfrac25\times300=120$ km. Avant la seconde pause :
$\dfrac7{15}\times300=140$ km.\\
4b) Fraction déjà parcourue : $\dfrac25+\dfrac7{15}=\dfrac6{15}+\dfrac7{15}=\dfrac{13}{15}$ ;
fraction restante : $\dfrac2{15}$, soit $\dfrac2{15}\times300=40$ km.

\corrige{du sujet 3}
1) $\dfrac{130}{195}=\dfrac{130\div65}{195\div65}=\dfrac23$.\\
2) Dénominateur commun $45$ : $-\dfrac45=-\dfrac{36}{45}$, $-\dfrac79=-\dfrac{35}{45}$. Comme
$-36<-35$, $-\dfrac45<-\dfrac79$.\\
3) $\dfrac54-\dfrac13=\dfrac{15}{12}-\dfrac4{12}=\dfrac{11}{12}$ ; puis
$\dfrac{11}{12}\times2=\dfrac{11}6$ ; enfin $E=\dfrac{11}6-\dfrac38=\dfrac{44}{24}-\dfrac9{24}=\dfrac{35}{24}$.\\
4a) Banque : $\dfrac25\times60\,000=24\,000$ FCFA ; il reste $60\,000-24\,000=36\,000$ FCFA.
Fournitures : $\dfrac14\times36\,000=9\,000$ FCFA.\\
4b) Il reste $36\,000-9\,000=27\,000$ FCFA ; part de la sœur : $\dfrac13\times27\,000=9\,000$
FCFA.\\
4c) Somme finale restant à Aïcha : $27\,000-9\,000=18\,000$ FCFA.

\end{multicols}
\exercicesBanner
